%==============================================================================
% Chapitre 32 - Précession
%==============================================================================

\section{Introduction}

Dans le chapitre précédent, nous avons vu qu'un projectile en rotation se
comporte comme un gyroscope.

Cette propriété lui confère une stabilité importante.

Cependant, un gyroscope ne réagit pas aux forces extérieures de manière
intuitive.

Lorsqu'un moment est appliqué, le mouvement observé apparaît souvent dans une
direction différente de celle attendue.

Ce phénomène est appelé précession.

\begin{aretenir}

La précession est le mouvement lent de rotation de l'axe d'un projectile sous
l'effet d'un moment extérieur.

\end{aretenir}

\begin{figure}[htbp]
\centering
\begin{tikzpicture}

\draw[dashed] (0,-2) -- (0,2);

\draw[thick] (0,0) -- (2,1.5);

\draw[dashed] (0,0) circle (1);

\draw[->] (1,0) arc (0:60:1);

\node at (2.5,1.7) {Axe du projectile};

\end{tikzpicture}
\caption{Principe de la précession.}
\label{fig:precession}
\end{figure}

\section{Exemple de la toupie}

Une toupie constitue l'exemple le plus connu.

Lorsque la gravité agit sur elle :

\begin{itemize}
    \item elle ne tombe pas immédiatement ;
    \item son axe tourne progressivement ;
    \item elle décrit un cône dans l'espace.
\end{itemize}

Ce mouvement correspond à la précession.

\section{Pourquoi est-ce contre-intuitif ?}

Pour un objet immobile :

\begin{itemize}
    \item une force produit une accélération dans la direction de cette force.
\end{itemize}

Pour un objet en rotation rapide :

\begin{itemize}
    \item la réponse est décalée ;
    \item le mouvement apparaît dans une autre direction.
\end{itemize}

C'est une conséquence directe de la conservation du moment cinétique.

\section{Application aux projectiles}

Pendant son vol, un projectile subit constamment :

\begin{itemize}
    \item des forces aérodynamiques ;
    \item des moments aérodynamiques ;
    \item des perturbations diverses.
\end{itemize}

Ces actions tendent à modifier son orientation.

En raison de sa rotation, la réponse prend la forme d'une précession.

\section{Incidence et précession}

Supposons qu'un projectile acquière une faible incidence.

L'écoulement d'air génère alors un moment aérodynamique.

Au lieu de simplement pivoter dans la direction attendue, le projectile
commence à précessionner.

Son axe effectue un mouvement lent autour de la direction moyenne du vol.

\section{Mouvement conique}

La précession peut être visualisée comme un mouvement conique.

L'extrémité de l'axe du projectile décrit une trajectoire proche d'un cercle.

Vu depuis l'avant, l'axe semble tourner autour de la trajectoire moyenne.

\section{Précession lente et précession rapide}

Les modèles avancés distinguent plusieurs composantes de précession.

On rencontre notamment :

\begin{itemize}
    \item une précession lente ;
    \item une précession rapide.
\end{itemize}

Ces composantes dépendent des caractéristiques du projectile et de sa
rotation.

\section{Précession et stabilité}

La précession n'est pas nécessairement un signe d'instabilité.

Au contraire :

\begin{itemize}
    \item elle fait partie du comportement normal d'un projectile stabilisé ;
    \item elle constitue une réponse naturelle aux moments appliqués.
\end{itemize}

\section{Amplitude de la précession}

L'amplitude dépend notamment :

\begin{itemize}
    \item de la vitesse de rotation ;
    \item de la stabilité gyroscopique ;
    \item de l'intensité des perturbations ;
    \item des propriétés aérodynamiques.
\end{itemize}

\section{Effet de la vitesse de rotation}

Une rotation élevée tend généralement à réduire certaines variations
d'orientation.

Le projectile résiste davantage aux perturbations.

Cependant, la relation exacte reste plus complexe qu'une simple augmentation
de la vitesse de rotation.

\section{Précession aérodynamique}

Dans le cas des projectiles, la précession est principalement entretenue par
les forces aérodynamiques.

Ces forces évoluent continuellement au cours du vol.

Le phénomène est donc dynamique.

\section{Amortissement}

L'atmosphère agit également comme un mécanisme d'amortissement.

Certaines oscillations diminuent progressivement.

Le projectile tend alors vers un comportement plus stable.

\section{Influence sur la trajectoire}

La précession modifie légèrement l'orientation du projectile.

Elle peut contribuer à :

\begin{itemize}
    \item certaines dérives ;
    \item certains mouvements oscillatoires ;
    \item l'apparition d'effets secondaires.
\end{itemize}

\section{Lien avec le spin drift}

La précession joue un rôle important dans l'apparition du spin drift.

Ce phénomène sera étudié dans un chapitre spécifique.

\section{Précession et simulation}

Les modèles simples ignorent généralement la précession.

Les modèles avancés la reproduisent explicitement.

Elle constitue l'un des phénomènes caractéristiques des modèles :

\begin{itemize}
    \item 5DOF ;
    \item 6DOF.
\end{itemize}

\begin{simulation}

La précession est souvent utilisée pour vérifier la cohérence d'un modèle
balistique avancé.

\end{simulation}

\section{Vers la nutation}

La précession décrit un mouvement relativement lent et régulier.

Cependant, un projectile réel présente souvent une seconde oscillation,
plus rapide et généralement plus faible.

Cette oscillation est appelée nutation.

\section{À retenir}

\begin{aretenir}

\begin{itemize}
    \item La précession est une conséquence directe de l'effet gyroscopique.
    \item Elle correspond à une rotation lente de l'axe du projectile.
    \item Elle apparaît lorsqu'un moment agit sur un projectile en rotation.
    \item Elle ne signifie pas nécessairement que le projectile est instable.
    \item Elle intervient dans plusieurs phénomènes balistiques avancés.
\end{itemize}

\end{aretenir}
